\documentclass[11pt,professionalfont,aspectratio=43]{beamer}

% Assets
\usepackage[czech]{babel}				% Jazyk
\usepackage[a-2u]{pdfx}				    % Kopírování z pdfka
\usepackage{tikz}						% Schémata automatů
\usetikzlibrary{shadows, shadows.blur, shapes.geometric, arrows, positioning, overlay-beamer-styles, patterns, shadings, shapes.callouts}
% \usetikzlibrary{arrows}
\usepackage{pgfplots}
% \pgfplotsset{compat=1.15}
% \usepackage{mathrsfs}
\usepackage[export]{adjustbox}

\usepackage[T1]{fontenc}
\usefonttheme[onlymath]{serif}
% \usepackage{newtxmath}
\usepackage{pxfonts}
\usepackage{slantsc}
% \usepackage{varwidth}

\usepackage[linesnumbered,ruled,vlined]{algorithm2e}                % Pseudokód algoritmů
\usepackage[utf8]{inputenc}
% \usepackage{textcomp}
\usepackage{hyperref}
\usepackage{fancybox}

%%% FONTS
% \usefonttheme{serif}
% \usepackage{lmodern}
\usepackage{helvet}

\hypersetup{%
    pdfencoding=auto,
    pdfauthor={\insertauthor},
    pdftitle={\insertsubtitle}
}
\usepackage{float}
\usepackage{csquotes}					% české uvozovky
\usepackage{enumerate}					% enumerate environment
\usepackage{indentfirst}
\usepackage{tcolorbox}                  % Fancy boxíky
\usepackage{mathtools}
\usepackage{pifont}
\usepackage{cancel}
\usepackage{subcaption}
\usepackage{soul}
\usepackage{makecell}                   % Zalomení v buňce tabulky
\usepackage{xcolor}
\usepackage{graphicx}
\usepackage{tabularx}
\usepackage{amsmath}
\usepackage{amsthm}
\usepackage{colortbl}                   % Barvení buněk v tabulce
\usepackage{multicol}                   % Rozdělení obsahu do sloucpů

\usepackage{listings}                   % Úryvky z kódu
\definecolor{maroon}{rgb}{0.5,0,0}
\definecolor{forestgreen}{rgb}{0.13, 0.55, 0.13}

% \lstdefinelanguage{XML}
% {
%   basicstyle=\ttfamily,
%   morestring=[s]{"}{"},
%   morecomment=[s]{?}{?},
%   morecomment=[s]{!--}{--},
%   commentstyle=\color{forestgreen},
%   moredelim=[s][\color{black}]{>}{<},
%   moredelim=[s][\color{red}]{\ }{=},
%   stringstyle=\color{blue},
%   identifierstyle=\color{maroon}
% }
% Colors

\def\maincolR{0} % červená složka (0–255)
\def\maincolG{83} % zelená složka
\def\maincolB{120} % modrá složka
\definecolor{maincolor}{HTML}{8B00B5}

% \pgfmathtruncatemacro{\authorfooterbgR}{1*\maincolR}
% \pgfmathtruncatemacro{\authorfooterbgG}{1*\maincolG}
% \pgfmathtruncatemacro{\authorfooterbgB}{1*\maincolB}
% \definecolor{authorfooterbg}{RGB}{\authorfooterbgR,\authorfooterbgG,\authorfooterbgB}

% \pgfmathtruncatemacro{\titlefooterbgR}{0.722*\maincolR}   
% \pgfmathtruncatemacro{\titlefooterbgG}{0.974*\maincolG}   
% \pgfmathtruncatemacro{\titlefooterbgB}{0.739*\maincolB}
% \definecolor{titlefooterbg}{RGB}{\titlefooterbgR,\titlefooterbgG,\titlefooterbgB}

% Beamer theme
\colorlet{themecolor}{maincolor}
\colorlet{authorfooterbg}{maincolor}
\definecolor{titlefooterbg}{HTML}{9471B0}
\definecolor{titlefooterfg}{HTML}{000000}
\colorlet{titlecolorfg}{maincolor}
\definecolor{titleboxbg}{HTML}{E7E7E7}
\definecolor{datefooterbg}{HTML}{E3E3E3}
\colorlet{titleframeboxcolor}{maincolor}

% Background gradient
\definecolor{gradMiddle}{RGB}{212,248,255}
\definecolor{gradEnd}{RGB}{114,231,255}

% Block theme
\definecolor{blocktitlebg}{HTML}{F0DFA9}
\definecolor{blocktitlefg}{HTML}{957300}
\definecolor{blockbodybg}{HTML}{F5F0DF}

\definecolor{alertblockbodybg}{HTML}{FFE4E4}
\definecolor{alertblocktitlebg}{HTML}{ECA4A4}
\definecolor{alertblocktitlefg}{HTML}{B11313}

\definecolor{exampleblockbodybg}{HTML}{DAF0DA}
\definecolor{exampleblocktitlebg}{HTML}{BBE0BB}
\definecolor{exampleblocktitlefg}{HTML}{257A25}

\definecolor{mathfg}{HTML}{00185E}

\colorlet{bulletbg}{maincolor}
% \definecolor{bulletbg}{HTML}{}

% Beamer theme
\usetheme{Madrid}
\usecolortheme{seahorse}
% \useoutertheme[subsection=false]{miniframes}

\usecolortheme[named=themecolor]{structure}
\setbeamertemplate{frame numbering}[fraction]
\setbeamertemplate{navigation symbols}{}
\setbeamercovered{invisible}

% Header/footer
\setbeamercolor{author in head/foot}{bg=authorfooterbg,fg=white}
\setbeamercolor{title in head/foot}{bg=titlefooterbg,fg=titlefooterfg}
\setbeamercolor{date in head/foot}{bg=datefooterbg,fg=black}
% \setbeamercolor{title}{bg=datefooterbg}
% \setbeamercolor{title}{fg=titlecolorfg}
\setbeamerfont{title}{series=\bfseries}
\setbeamercolor{frametitle}{bg=datefooterbg}
% \setbeamercolor{section in head/foot}{bg=red,fg=cyan}
% \setbeamercolor{subsection in head/foot}{bg=blue,fg=yellow}

% Table of contents
\setbeamercolor{section in toc}{fg=black}
\setbeamercolor{subsection in toc}{fg=red}
\setbeamercolor{section number projected}{fg=black}

% Blocks
\setbeamertemplate{blocks}[rounded][shadow=true]

\setbeamercolor{block title}{bg=blocktitlebg, fg=blocktitlefg}
\setbeamercolor{block body}{parent=normal text,use=block title,bg=blockbodybg}

\setbeamercolor{block body alerted}{bg=alertblockbodybg}
\setbeamercolor{block title alerted}{bg=alertblocktitlebg,fg=alertblocktitlefg}

\setbeamercolor{block body example}{bg=exampleblockbodybg}
\setbeamercolor{block title example}{bg=exampleblocktitlebg,fg=exampleblocktitlefg}

% Math
\setbeamercolor{math text}{fg=mathfg}

% Itemize a enumerate
\setbeamertemplate{itemize items}[circle]
\setbeamertemplate{enumerate items}[circle]
\setbeamercolor{itemize item}{fg=bulletbg,bg=white}
\setbeamercolor{itemize subitem}{fg=bulletbg,bg=white}
\setbeamercolor{itemize subsubitem}{fg=bulletbg,bg=white}

\setbeamertemplate{enumerate items}[circle]
\setbeamercolor{item projected}{bg=bulletbg}

% Titlepage
\defbeamertemplate*{title page}{customized}[1][]
{
    \vbox{}%
    \vfill%
    \begin{centering}
        \begin{tikzpicture}
            \node[thick, rectangle, rounded corners=5mm, drop shadow={shadow xshift=0mm, shadow yshift=2mm, opacity=0.2, blur shadow, shadow blur steps=10}, inner sep=5mm, fill=titleboxbg, fill opacity=1] {%
                {\usebeamerfont{title}\textcolor{titlecolorfg}{\inserttitle}\par}%
            };
        \end{tikzpicture}
        \vskip1em\par%
        \ifx\insertsubtitle\@empty%
        \else%
            {\usebeamerfont{subtitle}\insertsubtitle\par}%
        \fi%
        \vskip1em\par%
        \ifx\insertauthor\@empty%
        \else%
            {\usebeamerfont{author}\insertauthor}\\%
            \texttt{\email}
        \fi
        \vskip0.3em\par%
        % \ifx\insertinstitute\@empty%
        % \else%
        %     {\usebeamerfont{institute}\insertinstitute\par}%
        % \fi%
        \vskip1em\par%
        \ifx\inserttitlegraphic\@empty%
        \else%
            \vskip1em\par%
            \inserttitlegraphic\par%
        \fi%
        \vskip1em\par%
        \ifx\insertdate\@empty%
        \else%
            {\usebeamerfont{date}\insertdate\par}%
        \fi%
    \end{centering}%
    \vfill%
}
% Enumerate
%\setlist[enumerate]{topsep=0pt,itemsep=-1ex,partopsep=1ex,parsep=1ex,label=(\arabic*)}

\MakeOuterQuote{"}

%%% Makra pro definice, věty, tvrzení, příklady, ... (vyžaduje baliček amsthm)

\theoremstyle{plain}
\newtheorem{thm}{Věta}[section]
\newtheorem{lem}[theorem]{Lemma}
\newtheorem{prop}[theorem]{Tvrzení}
\newtheorem{cor}[theorem]{Důsledek}
\newtheorem*{prop*}{Tvrzení}

\theoremstyle{definition}
\newtheorem{define}[theorem]{Definice}
% \newtheorem{example}[theorem]{Příklad}
% \newtheorem{remark}[theorem]{Poznámka}
% \newtheorem{convention}[theorem]{Úmluva}
% \newtheorem{denoting}[theorem]{Značení}

\newenvironment{defblock}[1][]{
  \begin{block}{\ifthenelse{\equal{#1}{}}
  {Definice}
  {Definice (#1)}}
}{\end{block}}
\newenvironment{thmblock}[1][]{
  \begin{block}{\ifthenelse{\equal{#1}{}}
  {Věta}
  {Věta (#1)}}
}{\end{block}}
\newenvironment{corblock}[1][]{
  \begin{block}{\ifthenelse{\equal{#1}{}}
  {Důsledek}
  {Důsledek (#1)}}
}{\end{block}}
\newenvironment{propblock}[1][]{
  \begin{block}{\ifthenelse{\equal{#1}{}}
  {Tvrzení}
  {Tvrzení (#1)}}
}{\end{block}}
\newenvironment{exblock}[1][]{
  \begin{exampleblock}{\ifthenelse{\equal{#1}{}}
  {Příklad}
  {Příklad (#1)}}
}{\end{exampleblock}}


% Colors
\definecolor{lightblue}{HTML}{009AD4}
\definecolor{lightgreen}{HTML}{68FF00}
\definecolor{darkgreen}{HTML}{00B600}
\definecolor{darkblue}{HTML}{0265b3}
\definecolor{darkred}{HTML}{DA0707}
\definecolor{lightred}{HTML}{FF5100}
\definecolor{orange}{HTML}{FFE000}
\definecolor{codeblue}{HTML}{007eed}
\definecolor{codegreen}{rgb}{0,0.6,0}
\definecolor{codegray}{rgb}{0.5,0.5,0.5}
\definecolor{codebeige}{HTML}{D4A000}
\definecolor{codepink}{HTML}{FF0055}
\definecolor{backcolour}{rgb}{0.95,0.95,0.92}

\definecolor{decproblemtitlefg}{HTML}{957300}
\definecolor{decproblembg}{HTML}{F5F0DF}

\definecolor{timportantboxcolor}{HTML}{A97100}
\definecolor{talertboxcolor}{HTML}{B70000}
\definecolor{tnoticeboxcolor}{HTML}{159C00}
\definecolor{tinfoboxcolor}{HTML}{004DFF}

\newcommand{\markred}[1]{\textcolor{lightred}{#1}}
\newcommand{\markdarkred}[1]{\textcolor{darkred}{#1}}
\newcommand{\markgreen}[1]{\textcolor{lightgreen}{#1}}
\newcommand{\markdarkgreen}[1]{\textcolor{darkgreen}{#1}}
\newcommand{\markorange}[1]{\textcolor{orange}{#1}}
\newcommand{\markblue}[1]{\textcolor{lightblue}{#1}}
\newcommand{\markdarkblue}[1]{\textcolor{darkblue}{#1}}

% Math mode settings
% \AtBeginDocument{
%     \setlength{\abovedisplayskip}{3pt}
%     \setlength{\belowdisplayskip}{-4pt}
% }

% Inline images
\newcommand{\inlineimgscale}{1.1}

% X and check mark
\newcommand{\cmark}{\markgreen{\ding{51}}}
\newcommand{\xmark}{\markred{\ding{55}}}

% Math
\newcommand{\R}{\mathbb{R}}
\newcommand{\C}{\mathbb{C}}
\newcommand{\N}{\mathbb{N}}
\newcommand{\Z}{\mathbb{Z}}
\newcommand{\Q}{\mathbb{Q}}

\newcommand{\T}[1]{#1^\top}                     % Tranpozice matice
\newcommand{\compl}[1]{#1^\complement}
\newcommand{\mapping}[3]{#1:#2 \rightarrow #3}
\newcommand{\mapto}[3]{#1:#2 \mapsto #3}
\newcommand{\set}[1]{\left\{#1\right\}}
\newcommand{\powset}{\mathcal{P}}
\newcommand{\code}[1]{\langle#1\rangle}
\DeclareMathOperator*{\argmin}{arg\,min}
\DeclareMathOperator*{\argmax}{arg\,max}

% Statistika a AI
\newcommand{\average}[1]{\overline{#1}}
\DeclareMathOperator{\median}{med}
\DeclareMathOperator{\Var}{var}
\DeclareMathOperator{\Cov}{cov}
\DeclareMathOperator{\MSE}{MSE}
\DeclareMathOperator{\SSE}{SSE}

\DeclareMathOperator{\TP}{TP}
\DeclareMathOperator{\TN}{TN}
\DeclareMathOperator{\FP}{FP}
\DeclareMathOperator{\FN}{FN}
\DeclareMathOperator{\Accuracy}{Acc}
\DeclareMathOperator{\Precision}{Prec}
\DeclareMathOperator{\Recall}{Rec}
\DeclareMathOperator{\Fscore}{F}
\DeclareMathOperator{\Gini}{Gini}

\DeclareMathOperator{\mspan}{span}
\DeclareMathOperator{\mrank}{rank}
\DeclareMathOperator{\tg}{tg}
\DeclareMathOperator{\id}{id}
\DeclareMathOperator{\dom}{dom}
\DeclareMathOperator{\rng}{rng}
\renewcommand{\emptyset}{\varnothing}
\newcommand{\binary}[1]{\ensuremath{\left\langle #1\right\rangle}_B}

% DFS in and out lists
\DeclareMathOperator{\inlist}{in}
\DeclareMathOperator{\outlist}{out}
% \DeclareMathOperator{\lowlist}{low}
\DeclareMathOperator{\key}{key}

% A* macro
\newcommand{\Astar}{$\textcolor{black}{\textup{A}^{*}}$}

% Complexity classes
\DeclareMathOperator{\classP}{P}
\DeclareMathOperator{\classNP}{NP}
\DeclareMathOperator{\classTime}{TIME}
\DeclareMathOperator{\classNTime}{NTIME}
\DeclareMathOperator{\classSpace}{SPACE}
\DeclareMathOperator{\classNSpace}{NSPACE}

\DeclareMathOperator{\SAT}{SAT}
\DeclareMathOperator{\THREESAT}{3-SAT}
\DeclareMathOperator{\UNSAT}{UNSAT}
\DeclareMathOperator{\NZMNA}{NzMna}
\DeclareMathOperator{\KLIKA}{Klika}
\DeclareMathOperator{\ROZVRH}{Rozvrh}
\DeclareMathOperator{\SUDOKU}{Sudoku}
\newcommand{\HALT}{\ensuremath{\textsc{Halt}}}
\newcommand{\FACTOR}{\ensuremath{\textsc{Factor}}}

\DeclareMathOperator{\regex}{RegE}

%%% Barevné boxíky (notice / info / important / alert)
%
% Šířka boxu se nezadává ručně -- box se přizpůsobí šířce svého obsahu.
% Obsah se vysází do prostředí varwidth, takže krátký text zůstane na
% jednom řádku; delší text se zalomí, nejvýše však na šířku \linewidth
% (tj. \textwidth mimo sloupce a minipage). Pokud je titulek širší než
% obsah, řídí šířku boxu titulek.

\newsavebox{\tboxcontentbox}
\newsavebox{\tboxtitlebox}
\newlength{\tboxwidth}
\newlength{\tboxpadding}

% Geometrie boxu (musí odpovídat stylu tautobox níže):
% 2 * (boxrule + boxsep + left)
\setlength{\tboxpadding}{\dimexpr 0.5mm+1mm+4mm\relax}
\setlength{\tboxpadding}{2\tboxpadding}

\tcbset{tautobox/.style={boxrule=0.5mm, boxsep=1mm, left=4mm, right=4mm,
                         fonttitle=\bfseries}}

% Začátek sběru obsahu do boxu
\newcommand{\tboxcollect}{%
  \begin{lrbox}{\tboxcontentbox}%
    \begin{varwidth}{\dimexpr\linewidth-\tboxpadding\relax}%
}

% Ukončení sběru a výpočet výsledné šířky
% #1 = titulek (pro měření; prázdný u variant bez titulku)
\newcommand{\tboxmeasure}[1]{%
    \end{varwidth}%
  \end{lrbox}%
  \sbox{\tboxtitlebox}{\bfseries #1}%
  \setlength{\tboxwidth}{\wd\tboxcontentbox}%
  \ifdim\wd\tboxtitlebox>\tboxwidth
    \setlength{\tboxwidth}{\wd\tboxtitlebox}%
  \fi
  \addtolength{\tboxwidth}{\tboxpadding}%
  \ifdim\tboxwidth>\linewidth
    \setlength{\tboxwidth}{\linewidth}%
  \fi
}

% Vysázení boxu s titulkem: #1 = barva, #2 = titulek
\newcommand{\tboxrenderbox}[2]{%
  \tboxmeasure{#2}%
  \begin{center}%
    \begin{tcolorbox}[tautobox, colback=#1!15, colframe=#1,
                      title={#2}, width=\tboxwidth]
      \usebox{\tboxcontentbox}%
    \end{tcolorbox}%
  \end{center}%
}

% Vysázení boxu bez titulku: #1 = barva
\newcommand{\tboxrenderframe}[1]{%
  \tboxmeasure{}%
  \begin{center}%
    \begin{tcolorbox}[tautobox, colback=#1!15, colframe=#1,
                      width=\tboxwidth]
      \usebox{\tboxcontentbox}%
    \end{tcolorbox}%
  \end{center}%
}

% Poznámky (notice)
\newenvironment{tnoticebox}[1]{%
  \def\tboxtitletext{#1}\tboxcollect
}{%
  \tboxrenderbox{tnoticeboxcolor}{\tboxtitletext}%
}

\newenvironment{tnoticeframe}{%
  \tboxcollect
}{%
  \tboxrenderframe{tnoticeboxcolor}%
}

% Informační bloky (info)
\newenvironment{tinfobox}[1]{%
  \def\tboxtitletext{#1}\tboxcollect
}{%
  \tboxrenderbox{tinfoboxcolor}{\tboxtitletext}%
}

\newenvironment{tinfoframe}{%
  \tboxcollect
}{%
  \tboxrenderframe{tinfoboxcolor}%
}

% Důležité bloky (important)
\newenvironment{timportantbox}[1]{%
  \def\tboxtitletext{#1}\tboxcollect
}{%
  \tboxrenderbox{timportantboxcolor}{\tboxtitletext}%
}

\newenvironment{timportantframe}{%
  \tboxcollect
}{%
  \tboxrenderframe{timportantboxcolor}%
}

% Varovné bloky (alert)
\newenvironment{talertbox}[1]{%
  \def\tboxtitletext{#1}\tboxcollect
}{%
  \tboxrenderbox{talertboxcolor}{\tboxtitletext}%
}

\newenvironment{talertframe}{%
  \tboxcollect
}{%
  \tboxrenderframe{talertboxcolor}%
}

% Decision problem
\newcommand{\decproblem}[3]{
    \begin{center}
        \begin{minipage}{.9\textwidth}
            \begin{table}[!htbp]
                \centering
                \bgroup
                \def\arraystretch{1.8}
                \begin{tabularx}{\textwidth}{|rX|}
                \hline
                \rowcolor{decproblembg} \multicolumn{2}{|c|}{\textcolor{decproblemtitlefg}{\textsc{#1}}} \\ \hline
                \markdarkred{\textbf{Vstup:}} & #2                \\ \hline
                \rowcolor{decproblembg} \markdarkred{\textbf{Otázka:}} & #3               \\ \hline
                \end{tabularx}
                \egroup
            \end{table}
        \end{minipage}
    \end{center}
}
\newcommand{\bigO}{\mathcal{O}}

% TODO
\newcommand{\todo}[1]{\textcolor{red}{(\noindent TODO: #1.)}}


\newcommand{\hlcode}[1]{\colorbox{red}{#1}}

% Algorithms
\numberwithin{algocf}{section}
\renewcommand{\algorithmcfname}{Algoritmus}
\SetKwInput{KwIn}{Vstup}
\SetKwInput{KwOut}{Výstup}
\SetKwProg{Function}{Funkce}{}{end}
\setlength{\algomargin}{25pt}
\renewcommand{\thealgocf}{}     % Vypnutí číslování algoritmů

% Definice stylů pro uzly
\tikzstyle{startstop} = [rectangle, rounded corners, minimum width=1.5cm, minimum height=0.5cm,text centered, draw=black, fill=red!30]
\tikzstyle{process}   = [rectangle, minimum width=1.5cm, minimum height=0.5cm, text centered, draw=black, fill=orange!30]
\tikzstyle{decision}  = [diamond, aspect=2, text centered, draw=black, fill=green!30]
\tikzstyle{arrow}     = [thick,->,>=stealth]

\tikzstyle{vertex}    = [thick, draw, circle, minimum size=4mm, inner sep=1pt, outer sep=1pt, align=center]
\tikzstyle{vertexplain}    = [minimum size=4mm, inner sep=1pt, outer sep=1pt, align=center]
\tikzstyle{verteximg}    = [draw, circle, minimum size=4mm, inner sep=0pt, outer sep=1pt, align=center, clip]
\tikzstyle{verteximgplain}    = [minimum size=4mm, inner sep=0pt, outer sep=1pt, align=center]
\tikzstyle{task}    = [draw, rounded corners=2pt, minimum width=14mm, minimum height=8mm, inner sep=2pt, align=center]

\tikzstyle{edge}    = [thick,>=stealth]
\tikzstyle{diredge}    = [thick, ->, >=stealth]

%% Automaty
\tikzstyle{state}=[%
    draw,
    circle,
    thick,
    minimum size=10mm,
    inner sep=0pt,
    outer sep=0pt
]

\tikzstyle{acceptstate}=[%
    state,
    double,
    double distance=1pt
]

\tikzstyle{transition}=[%
    ->,
    >=stealth,
    thick
]

\newcommand<>{\state}[4][]{%
    \node#5[state,#1] (#2) at #3 {#4};
}

\newcommand<>{\acceptingstate}[4][]{%
    \node#5[acceptstate,#1] (#2) at #3 {#4};
}

\newcommand<>{\transitionarrow}[5][]{%
    \path#6 (#2) edge[transition,#1] node[#3] {#4} (#5);
}

\newcommand<>{\initialstate}[6][]{%
    \node#7[state,#1] (#2) at #3 {#4};
    \draw#7[transition] #5 -- (#2.#6);
}

%%%%%

% Frames
\newcommand{\tableofcontentsframe}{
    \begin{frame}{Obsah}
        \tableofcontents
    \end{frame}
}

% \definecolor{boxcolor}{HTML}{00820a}
\newcommand{\titleframe}[1]{%
  \begin{frame}[plain]
    \begin{tikzpicture}[remember picture, overlay]
      \node at (current page.center)
        [draw=titleframeboxcolor, line width=1pt,
         inner xsep=10mm, inner ysep=4mm, rounded corners=2mm,
         text width=0.8\textwidth, align=center,
         execute at begin node=\setlength{\baselineskip}{2.2em}]
        {\Huge\markred{#1}\normalfont};
    \end{tikzpicture}
  \end{frame}
  \addtocounter{framenumber}{-1}%
}

\newenvironment{titleframeenv}[1]{
    \begin{frame}[plain]
        \begin{tcolorbox}[colback=white,boxrule=1pt]
            \begin{center}
                \Huge\markred{#1}\normalfont
            \end{center}
        \end{tcolorbox}
}{\addtocounter{framenumber}{-1}\end{frame}}

\newcommand{\icon}[1]{
  \includegraphics[height=\fontcharht\font`\B]{#1}
}

% Obrázek, který se vždy vejde na snímek. Adjustbox jej v případě potřeby
% zmenší tak, aby nepřesáhl šířku textu ani zadanou část výšky snímku; díky
% tomu sedí sazba ve všech používaných poměrech stran (4:3, 16:9 i 16:10).
% Volitelným argumentem je maximální výška (výchozí 0.68\textheight), povinným
% cesta k souboru s obrázkem (TikZ zdroj vkládaný přes \input).
\newcommand{\obrazek}[2][0.68\textheight]{%
    \begin{figure}
        \centering
        \begin{adjustbox}{max width=\linewidth, max totalheight=#1, center}
            \input{#2}
        \end{adjustbox}
    \end{figure}%
}

% Totéž pro rastrové či vektorové obrázky vkládané přes \includegraphics.
\newcommand{\obrazekgrf}[2][0.68\textheight]{%
    \begin{figure}
        \centering
        \includegraphics[width=\linewidth, max totalheight=#1, center]{#2}
    \end{figure}%
}

% Závěrečný snímek se seznamem zdrojů. Prvním (volitelným) argumentem je název
% sekce i snímku, druhým čárkou oddělený seznam klíčů z ../assets/literatura.bib.
% Citace se sázejí podle ČSN ISO 690, v pořadí, v jakém jsou klíče uvedeny.
\newcommand{\zdrojeframe}[2][Zdroje]{%
    \section{#1}
    \nocite{#2}%
    \begin{frame}[allowframebreaks]{#1}
        \printbibliography[heading=none]
    \end{frame}
}
\input{../assets/listing.tex}

\def\presentationtitle{Lineární regrese}

\begin{document}

\title[\presentationtitle]{\textsc{Teoretická informatika}}
\subtitle{\presentationtitle}
\author{David Weber}
\def\email{weber3@spsejecna.cz}
\institute{SPŠE Ječná}
\date{%
    Rok 
    \pgfmathsetmacro{\firstyear}{ifthenelse(\month<9, \year-1, \year)}%
    \pgfmathsetmacro{\secondyear}{ifthenelse(\month<9, \year, \year+1)}%
    \pgfmathtruncatemacro{\secondyearlasttwo}{mod(\secondyear,100)}%
    \pgfmathprintnumber[fixed,precision=0,zerofill,1000 sep={}]{\firstyear}/%
    \pgfmathprintnumber[fixed,precision=0,zerofill,1000 sep={}]{\secondyearlasttwo}%
}
\titlegraphic{\includegraphics[width=0.25\textwidth]{../images/SPSE-Jecna_Logo.pdf}}

\begin{frame}[plain]
    \titlepage
\end{frame}

% =================================================================
\section{Lineární regrese}
% =================================================================

\begin{titleframeenv}{Lineární regrese}
    \obrazek[0.4\textheight]{images/metoda-nejmensich-ctvercu.tex}
\end{titleframeenv}

\begin{frame}{Motivace: z~dat k~předpovědi}

    \visible<2->{\begin{columns}[T]
        \begin{column}{0.46\textwidth}
            \begin{itemize}
                \item Máme dvojice hodnot
                \[
                    (x_1,y_1),\ldots,(x_n,y_n).
                \]
                \visible<3->{\item Z~hodnoty $x$ chceme odhadnout odpovídající hodnotu $y$.}
                \visible<4->{\item Bodový graf může naznačovat přibližně lineární vztah.}
            \end{itemize}
        \end{column}

        \begin{column}{0.50\textwidth}\visible<5->{
    \obrazek[0.4\textheight]{images/motivace-z-dat-k-predpovedi.tex}
        }\end{column}
    \end{columns}}

    \visible<6->{\begin{timportantframe}
        \begin{center}
            Jak najít přímku, která body vystihuje nejlépe?
        \end{center}
    \end{timportantframe}}
\end{frame}

\subsection{Co je lineární regrese?}

\begin{frame}{Co je lineární regrese?}

    \visible<2->{Lineární regrese hledá lineární funkci
    \[
        f(x)=ax+b,
    \]
    která \markdarkred{co nejlépe popisuje vztah} mezi statistickými soubory
    \[
        \mathbf X=(x_1,\ldots,x_n),
        \qquad
        \mathbf Y=(y_1,\ldots,y_n).
    \]}


    \visible<3->{\begin{itemize}
        \item Hodnota $f(x_i)$ je předpověď modelu pro vstup $x_i$.
        \visible<4->{\item Koeficient $a$ určuje směrnici přímky.}
        \visible<5->{\item Koeficient $b$ určuje její posunutí.}
    \end{itemize}}

    \visible<6->{\begin{tinfoframe}
        Cílem není najít přímku procházející všemi body, ale přímku s~co nejmenší celkovou chybou.
    \end{tinfoframe}}
\end{frame}

\subsection{Predikce a~rezidua}

\begin{frame}{Predikce a~rezidua}

    \visible<2->{Pro každé pozorování budeme značit
    \[
        \hat y_i=f(x_i)
    \]
    jako předpovězenou hodnotu a
    \[
        e_i=y_i-\hat y_i
    \]
    nazýváme \markdarkred{reziduum}.}

    \visible<3->{\vspace{0.2cm}

    \begin{columns}[T]
        \begin{column}{0.45\textwidth}
            \begin{itemize}
                \item $e_i>0$: bod leží nad přímkou,
                \item $e_i<0$: bod leží pod přímkou,
                \item $e_i=0$: bod leží přesně na přímce.
            \end{itemize}
        \end{column}

        \begin{column}{0.50\textwidth}\visible<4->{
    \obrazek[0.4\textheight]{images/predikce-a-rezidua.tex}
        }\end{column}
    \end{columns}}
\end{frame}

\subsection{Měření chyby}

\begin{frame}{Jak měřit celkovou chybu modelu?}
    \begin{itemize}
        \item Pouhý součet reziduí není vhodný:
        \[
            \sum_{i=1}^n e_i
        \]
        \visible<2->{protože kladné a~záporné chyby se mohou vyrušit.
        \item Proto používáme čtverce reziduí:
        \[
            e_i^2=(y_i-f(x_i))^2.
        \]}
    \end{itemize}
\end{frame}

\subsection{Součet čtvercových chyb a~střední kvadratická chyba}

\begin{frame}{SSE a~MSE}
    \begin{defblock}
        Nechť jsou dány statistické soubory $\mathbf X=(x_1,\ldots,x_n)$ a~$\mathbf Y=(y_1,\ldots,y_n)$ a~funkce $\mapping{f}{\R}{\R}$. Pak
        \begin{itemize}
            \item \markdarkred{součet čtvercových chyb} (sum of squared errors) funkce $f$ definujeme jako
            \[
                \SSE_{\mathbf X,\mathbf Y}(f)
                =
                \sum_{i=1}^n
                \left(y_i-f(x_i)\right)^2
                =
                \sum_{i=1}^n e_i^{2}
                .
            \]
            \item \markdarkred{střední kvadratickou chybu} (mean squared error) funkce $f$ definujeme jako
            \[
                \MSE_{\mathbf X,\mathbf Y}(f)
                =
                \frac1n\SSE_{\mathbf X,\mathbf Y}(f).
            \]
        \end{itemize}
    \end{defblock}
\end{frame}

\subsection{Metoda nejmenších čtverců}

\begin{frame}{Metoda nejmenších čtverců}
    \small

    \visible<2->{\begin{itemize}
        \item Metoda nejmenších čtverců vybírá takovou funkci $f$, pro kterou je hodnota
        \[
            \sum_{i=1}^n
            \left(y_i-f(x_i)\right)^2
        \]
        nejmenší.
        \visible<3->{\item Pro lineární regresi tedy hledáme dvojici
        \[
            (a,b)=\argmin_{(a,b)\in\mathbb R^{2}}
            \sum_{i=1}^n
            (y_i-ax_i-b)^2.
        \]}
        \visible<4->{\item Takovou optimální přímku budeme nazývat \markdarkred{regresní přímka souborů $\mathbf{X}$ a~$\mathbf{Y}$}, nebo též jejich \markdarkred{lineární model}.}
    \end{itemize}}

    \visible<5->{\begin{tinfoframe}
        Minimalizace $\SSE$ a~$\MSE$ vede ke stejnému modelu, protože se liší pouze násobkem $\frac1n$. Proto budeme dále používat pro zjednodušení situace $\SSE$.
    \end{tinfoframe}}
\end{frame}


\begin{frame}{Metoda nejmenších čtverců}
    \obrazek[0.4\textheight]{images/metoda-nejmensich-ctvercu.tex}

    \begin{timportantframe}
        Hledáme přímku, pro kterou je součet obsahů těchto čtverců nejmenší.
    \end{timportantframe}
\end{frame}


\begin{frame}{Souvislost s~aritmetickým průměrem}

    \visible<2->{V~prezentaci o~statistice jsme ukázali, že aritmetický průměr minimalizuje součet čtvercových odchylek:
    \[
        \average{\mathbf{X}}
        =
        \argmin_{c\in\mathbb R}
        \overbrace{\sum_{i=1}^n (x_i-c)^2}^{\SSE_{\mathbf X,\mathbf X}(x\mapsto c)}.
    \]}



    \visible<3->{\begin{tinfoframe}
        Aritmetický průměr je tedy nejlepší konstantní odhad hodnot souboru $\mathbf X$ ve~smyslu čtvercové chyby.
    \end{tinfoframe}}



    \visible<4->{Lineární regrese tuto myšlenku zobecňuje:
    \begin{itemize}
        \item místo jedné konstanty hledáme přímku,
        \item místo odchylek od čísla měříme odchylky od hodnot $f(x_i)$.
    \end{itemize}}
\end{frame}


\begin{frame}{Nejlepší konstantní model}

    \visible<2->{Nejprve uvažujme pouze konstantní model
    \[
        f(x)=b.
    \]
    pro soubory $\mathbf{X},\mathbf{Y}$. Chceme minimalizovat výraz
    \[
        \SSE_{\mathbf{X},\mathbf{Y}}(f)=\sum_{i=1}^n
        (y_i-b)^2.
    \]}



    \visible<3->{Podle tvrzení o~aritmetickém průměru platí
    \[
        b=\average{\mathbf{Y}},\quad \text{a tedy}\;f(x)=\average{\mathbf{Y}}.
    \]}

    \visible<4->{\begin{tinfoframe}
        \begin{center}
            Tento model později porovnáme s~regresní přímkou.
        \end{center}
    \end{tinfoframe}}
\end{frame}

\subsection{Hledání nejlepší přímky}

\begin{frame}{Hledání nejlepší přímky}

    \visible<2->{Pro lineární model
    \[
        f_{a,b}(x)=ax+b
    \]
    se soubory $\mathbf{X},\mathbf{Y}$ definujeme chybovou funkci
    \[
        E(a,b)=\SSE_{\mathbf{X},\mathbf{Y}}(f_{a,b})=\sum_{i=1}^n(y_i-ax_i-b)^2.
    \]}

    \visible<3->{\begin{timportantframe}
        Cílem je najít dvojici $(a,b)$ takovou, aby hodnota $E(a,b)$ byla nejmenší.
    \end{timportantframe}}

    \visible<4->{Výpočet rozdělíme na dva kroky:}

    \visible<5->{\begin{enumerate}
        \item najdeme nejlepší $b$ pro pevně zvolené $a$,
        \visible<6->{\item potom najdeme nejlepší směrnici $a$.}
    \end{enumerate}}
\end{frame}


\begin{frame}{Nejlepší hodnota $b$ pro pevné $a$}

    \visible<2->{\begin{itemize}
        \item Pro pevné $a$ můžeme psát
        \[
            E(a,b)
            =
            \sum_{i=1}^n(y_i-ax_i-b)^2
            =
            \sum_{i=1}^n
            \left((y_i-ax_i)-b\right)^2.
        \]
        \visible<3->{\item To je opět úloha hledání nejlepší konstanty pro soubor hodnot
        \[
            \mathbf{Y}-a\mathbf{X}=(y_1-ax_1,\dots,y_n-ax_n).
        \]}
        \visible<4->{\item Proto dostáváme
        \[
            b=\average{\mathbf{Y}-a\mathbf{X}}=\average{\mathbf{Y}}-a\average{\mathbf{X}}.
        \]}
    \end{itemize}}
    
\end{frame}


\begin{frame}{Přímka procházející bodem průměrů}

    \visible<2->{\begin{itemize}
        \item Po dosazení
        \[
            b=\average{\mathbf{Y}}-a\average{\mathbf{X}}
        \]
        \visible<3->{lze model zapsat ve tvaru
        \[
            f(x)=\average{\mathbf{Y}}+a(x-\average{\mathbf{X}}).
        \]}
        \visible<4->{\item Odtud ihned plyne
        \[
            f(\average{\mathbf{X}})=\average{\mathbf{Y}}.
        \]}
    \end{itemize}}

    \visible<5->{\begin{tnoticeframe}
        \begin{center}
            Každý "kandidát" na~nejlepší přímku prochází bodem $(\average{\mathbf{X}},\average{\mathbf{Y}})$.
        \end{center}
    \end{tnoticeframe}}

\end{frame}


\begin{frame}{Přímky procházející bodem $(\average{\mathbf{X}},\average{\mathbf{Y}})$}
    \obrazek[0.4\textheight]{images/primky-prochazejici-bodem-x-y.tex}

    \begin{tinfoframe}
        Zbývá zvolit směrnici $a$ tak, aby výsledná přímka měla nejmenší chybu $E(a,b)$.
    \end{tinfoframe}
\end{frame}


\begin{frame}{Odvození směrnice $a$}
    \small

    \visible<2->{Po dosazení
    \[
        b=\average{\mathbf{Y}}-a\average{\mathbf{X}}
    \]
    dostaneme funkci jedné proměnné:
    \[
        E(a)=E(a,\average{\mathbf{Y}}-a\average{\mathbf{X}})
        =
        \sum_{i=1}^n
        \left(y_i-ax_i-(\average{\mathbf{Y}}-a\average{\mathbf{X}})\right)^2
        =
        \sum_{i=1}^n
        \left(
        (y_i-\average{\mathbf{Y}})
        -
        a(x_i-\average{\mathbf{X}})
        \right)^2.
    \]}



    \visible<3->{Po roznásobení:
    \[
        E(a)
        =
        \sum_{i=1}^n
        (y_i-\average{\mathbf{Y}})^2 -\boxed{a}\cdot 2\sum_{i=1}^n
        (x_i-\average{\mathbf{X}})(y_i-\average{\mathbf{Y}})+\boxed{a^2}\sum_{i=1}^n (x_i-\average{\mathbf{X}})^2.
    \]}
\end{frame}


\begin{frame}{Vzorec pro směrnici regresní přímky}

    \visible<2->{\begin{itemize}
        \item Funkce $E(a)$ je kvadratická funkce v~proměnné $a$.
        \visible<3->{\item Minimum nastává pro
        \[
            a=
            \frac{
            \sum_{i=1}^n
            (x_i-\average{\mathbf{X}})(y_i-\average{\mathbf{Y}})
            }{
            \sum_{i=1}^n
            (x_i-\average{\mathbf{X}})^2
            }
            =
            \frac{
            \frac1n\sum_{i=1}^n
            (x_i-\average{\mathbf{X}})(y_i-\average{\mathbf{Y}})
            }{
            \frac1n\sum_{i=1}^n
            (x_i-\average{\mathbf{X}})^2
            }
            =
            \boxed{\frac{\Cov(\mathbf{X},\mathbf{Y})}{\Var(\mathbf{X})}}.
        \]}
        \visible<4->{\item Jakmile známe $a$, dopočítáme
        \[
            \boxed{
            b=\average{\mathbf{Y}}-a\average{\mathbf{X}}
            }.
        \]}
    \end{itemize}}

    \visible<5->{\begin{talertframe}
        \begin{center}
            \textbf{Podmínka:} hodnoty $x_i$ nesmějí být všechny stejné.
        \end{center}
    \end{talertframe}}
\end{frame}

\subsection{Tvrzení o~nejlepší regresní přímce}

\begin{frame}{Tvrzení o~nejlepší regresní přímce}
    \begin{propblock}
        Nechť jsou dány statistické soubory $\mathbf{X},\mathbf{Y}$. Potom \markdarkred{existuje právě jedna} lineární funkce
        \[
            f(x)=ax+b,
        \]
        která minimalizuje součet čtvercových chyb, tj. 
        \[
            f=\argmin_{\ell\;\text{lineární}}\;\SSE_{\mathbf{X},\mathbf{Y}}(\ell)
        \]
        a~její koeficienty jsou
        \[
            a=\frac{\Cov(\mathbf{X},\mathbf{Y})}{\Var(\mathbf{X})},
            \qquad
            b=\average{\mathbf{Y}}-a\average{\mathbf{X}}.
        \]
    \end{propblock}
\end{frame}


\begin{frame}{Geometrická interpretace výsledku}

    \visible<2->{\begin{itemize}
        \item Regresní přímka prochází bodem
        \[
            (\average{\mathbf{X}},\average{\mathbf{Y}}).
        \]
        \visible<3->{\item Minimalizuje součet čtverců svislých odchylek bodů od přímky.}
    \end{itemize}}

    \visible<4->{\begin{talertframe}
        Směr regrese je podstatný:
        \[
            \text{regrese }\mathbf{Y}\text{ podle }\mathbf{X}
            \quad\neq\quad
            \text{regrese }\mathbf{X}\text{ podle }\mathbf{Y}.
        \]
        Výsledné přímky nejsou stejné.
    \end{talertframe}}
\end{frame}

\subsection{Směrnice regresní přímky a~korelace}

\begin{frame}{Směrnice regresní přímky a~korelace}

    \visible<2->{\begin{itemize}
        \item Korelační koeficient je
        \[
            \rho_{\mathbf{X}\mathbf{Y}}
            =
            \frac{\Cov(\mathbf X,\mathbf Y)}
            {\sigma_\mathbf{X}\sigma_\mathbf{Y}}.
        \]
        \visible<3->{\item Proto $\Cov(\mathbf{X},\mathbf{Y})=\rho_{\mathbf{X}\mathbf{Y}}\sigma_\mathbf{X}\sigma_\mathbf{Y}$.}
        \visible<4->{\item Po dosazení do vzorce pro $a$ dostaneme
        \[
            a=\frac{\Cov(\mathbf X,\mathbf Y)}{\Var(\mathbf X)}=\frac{\rho_{\mathbf{X}\mathbf{Y}}\sigma_\mathbf{X}\sigma_\mathbf{Y}}{\sigma^{2}_\mathbf{X}}=\rho_{\mathbf{X}\mathbf{Y}}\frac{\sigma_\mathbf{Y}}{\sigma_\mathbf{X}}.
        \]}
        \visible<5->{\item Regresní přímku můžeme tedy vyjádřit jako
        \[
            \boxed{
            f(x)=
            \average{\mathbf{Y}}+
            \rho_{\mathbf{X}\mathbf{Y}}
            \frac{\sigma_\mathbf{Y}}{\sigma_\mathbf{X}}
            (x-\average{\mathbf{X}})
            }.
        \]}
    \end{itemize}}
    
\end{frame}


\begin{frame}{Co ze vztahu se směrnicí plyne?}

    \visible<2->{Ze vztahu
    \[
        a=
        \rho_{\mathbf{X}\mathbf{Y}}
        \frac{\sigma_\mathbf{Y}}{\sigma_\mathbf{X}}
    \]
    plyne:}

    \visible<3->{\begin{itemize}
        \item znaménko směrnice je stejné jako znaménko korelace,
        \visible<4->{\item pokud $\rho_{\mathbf{X}\mathbf{Y}}>0$, přímka je rostoucí,}
        \visible<5->{\item pokud $\rho_{\mathbf{X}\mathbf{Y}}<0$, přímka je klesající,}
        \visible<6->{\item pokud $\rho_{\mathbf{X}\mathbf{Y}}=0$, potom $a=0$.}
    \end{itemize}}

    \visible<6->{\begin{corblock}
        Jestliže $\rho_{\mathbf{X}\mathbf{Y}}=0$, nejlepší lineární model je $f(x)=\average{\mathbf{Y}}$.
    \end{corblock}}
\end{frame}

\subsection{Korelace jako směrnice po standardizaci}

\begin{frame}{Korelace jako směrnice po standardizaci}

    \visible<2->{Definujme standardizované hodnoty souborů $\mathbf{X}$ a~$\mathbf{Y}$
    \[
        u_i=
        \frac{x_i-\average{\mathbf{X}}}{\sigma_\mathbf{X}},
        v_i=
        \frac{y_i-\average{\mathbf{Y}}}{\sigma_\mathbf{Y}}
        \quad
        \text{a}
        \quad
        \mathbf{U}=(u_1,\dots,u_n),\quad\mathbf{V}=(v_1,\dots,v_n),
    \]}
    \visible<3->{Pro standardizované soubory $\mathbf{U}$ a~$\mathbf{V}$ má regresní přímka tvar
    \[
        \boxed{
        v=\rho_{\mathbf{X}\mathbf{Y}}u
        }.
    \]}

    \visible<4->{\begin{timportantframe}
        Korelační koeficient je směrnicí regresní přímky mezi standardizovanými hodnotami.
    \end{timportantframe}}

    \visible<5->{\begin{itemize}
        \item Korelační koeficient nemá jednotky.
        \visible<6->{\item Změna jednotek nezmění korelační koeficient.}
    \end{itemize}}
\end{frame}

\subsection{Minimální chyba a~korelace}

\begin{frame}{Minimální chyba a~korelace}

    \visible<2->{\begin{propblock}[Bez důkazu]
        Pro soubory $\mathbf{X},\mathbf{Y}$ a~jejich regresní přímku $f$ platí
        \[
            \MSE_{\mathbf{X},\mathbf{Y}}(f)
            =
            \sigma_\mathbf{Y}^2(1-\rho_{\mathbf{X}\mathbf{Y}}^2).
        \]
    \end{propblock}}



    \visible<3->{\begin{itemize}
        \item Pokud $|\rho_{\mathbf{X}\mathbf{Y}}|=1$, potom
        \[
            \MSE_{\mathbf{X},\mathbf{Y}}(f)=0.
        \]
        Všechny body leží přesně na jedné přímce.

        \visible<4->{\item Pokud $\rho_{\mathbf{X}\mathbf{Y}}=0$, potom
        \[
            \MSE_{\mathbf{X},\mathbf{Y}}(f)=\sigma_\mathbf{Y}^2=\Var(\mathbf{Y}).
        \]
        Nejlepší lineární model se nezlepší oproti konstantnímu modelu.}
    \end{itemize}}
\end{frame}

% =================================================================
\section{Koeficient determinace}
% =================================================================

\begin{titleframeenv}{Koeficient determinace}
    \obrazek[0.4\textheight]{images/koeficient-determinace.tex}
\end{titleframeenv}

\subsection{Definice}

\begin{frame}{Koeficient determinace}
    \small

    \visible<2->{Chybu regresní přímky porovnáme s~chybou konstantního modelu
    \[
        f(x)=\average{\mathbf{Y}}.
    \]}

    \visible<3->{\begin{defblock}[Koeficient determinace]
        Nechť jsou dány soubory $\mathbf X=(x_1,\ldots,x_n)$ a~$\mathbf Y=(y_1,\ldots,y_n)$, kde $\sigma_\mathbf{Y}>0$ a~dále $f$ je jejich regresní přímka. \markdarkred{Koeficient determinace} modelu $f$ definujeme jako
        \[
            R^2=1-\frac{\MSE_{\mathbf X,\mathbf Y}(f)}{\MSE_{\mathbf X,\mathbf Y}(x\mapsto\average{\mathbf{Y}})}=1-\frac{\MSE_{\mathbf X,\mathbf Y}(f)}{\Var(\mathbf{Y})}.
        \]
    \end{defblock}}

    \visible<4->{Pro jednoduchou lineární regresi platí
    \[
        R^2=\rho_{\mathbf{X}\mathbf{Y}}^2.
    \]}
\end{frame}

\begin{frame}{Koeficient determinace}
    \obrazek[0.4\textheight]{images/koeficient-determinace.tex}
\end{frame}

\subsection{Příklad výpočtu}

\begin{frame}{Příklad}

    \visible<2->{Uvažujme statistické soubory
    \[
        \mathbf X=(1,2,3,4,5),
        \qquad
        \mathbf Y=(2,3,5,4,6).
    \]}

    \visible<3->{Dvojice $(x_i,y_i)$ představují jednotlivá pozorování:
    \[
        \begin{array}{c|ccccc}
            i   & 1 & 2 & 3 & 4 & 5 \\ \hline
            x_i & 1 & 2 & 3 & 4 & 5 \\
            y_i & 2 & 3 & 5 & 4 & 6
        \end{array}
    \]}

    \visible<4->{\obrazek[0.28\textheight]{images/priklad.tex}}
\end{frame}

\begin{frame}{Příklad}{Výpočet regresní přímky}
    \small

    \visible<2->{Nejprve spočítáme aritmetické průměry:
    \[
        \average{\mathbf{X}}
        =
        \frac{1+2+3+4+5}{5}
        =
        3,\quad
        \average{\mathbf{Y}}
        =
        \frac{2+3+5+4+6}{5}
        =
        4.
    \]}

    \visible<3->{Směrnice regresní přímky je
    \[
        a
        =
        \frac{
            \sum_{i=1}^5
            (x_i-\average{\mathbf{X}})(y_i-\average{\mathbf{Y}})
        }{
            \sum_{i=1}^5
            (x_i-\average{\mathbf{X}})^2
        }
        =\frac{
            4+1+0+0+4
        }{
            4+1+0+1+4
        }
        =
        \frac9{10}
        =
        0{,}9.
    \]}

    \visible<4->{Absolutní člen je
    \[
        b
        =
        \average{\mathbf{Y}}-a\average{\mathbf{X}}
        =
        4-0{,}9\cdot 3
        =
        1{,}3.
    \]}

    \visible<5->{\begin{timportantframe}
        \begin{center}
            Hledaná regresní přímka je $f(x)=0{,}9x+1{,}3$.
        \end{center}
    \end{timportantframe}}
\end{frame}

% 3
\begin{frame}{Příklad}{Předpovědi a~rezidua}
    \small

    \visible<2->{Pro regresní model
    \[
        f(x)=0{,}9x+1{,}3
    \]
    vypočítáme předpovědi $\hat y_i=f(x_i)$ a~rezidua $e_i=y_i-\hat y_i$.}

    \visible<3->{\[
        \begin{array}{c|ccccc}
            i
                & 1 & 2 & 3 & 4 & 5 \\ \hline
            x_i
                & 1 & 2 & 3 & 4 & 5 \\
            y_i
                & 2 & 3 & 5 & 4 & 6 \\
            \hat y_i
                & 2{,}2 & 3{,}1 & 4{,}0 & 4{,}9 & 5{,}8 \\
            e_i
                & -0{,}2 & -0{,}1 & 1{,}0 & -0{,}9 & 0{,}2
        \end{array}
    \]}

    \visible<4->{\begin{align*}
        \SSE_{\mathbf{X},\mathbf{Y}}(f)&=\sum_{i=1}^5 e_i^2=0{,}04+0{,}01+1+0{,}81+0{,}04=1{,}9.\\
        \MSE_{\mathbf{X},\mathbf{Y}}(f)&=\frac{1{,}9}{5}=0{,}38.
    \end{align*}}
\end{frame}

\begin{frame}{Příklad}{Předpovědi a~rezidua}
    \obrazek[0.4\textheight]{images/priklad-predpovedi-a-rezidua.tex}
\end{frame}

% 4
\begin{frame}{Příklad}{Koeficient determinace}

    \visible<2->{Nejlepší konstantní model je
    \[
        c(x)=\average{\mathbf{Y}}=4.
    \]}

    \visible<3->{Jeho chyba je
    \[
    \begin{aligned}
        \MSE_{\mathbf X,\mathbf Y}(c)
        &=
        \frac15
        \sum_{i=1}^5
        (y_i-\average{\mathbf{Y}})^2=
        \frac15
        \left(
            (-2)^2+(-1)^2+1^2+0^2+2^2
        \right)=2.
    \end{aligned}
    \]}

    \visible<4->{Pro regresní přímku jsme vypočítali $\MSE_{\mathbf X,\mathbf Y}(f)=0{,}38$.}

    \visible<5->{Proto
    \[
        R^2
        =
        1-
        \frac{
            \MSE_{\mathbf X,\mathbf Y}(f)
        }{
            \MSE_{\mathbf X,\mathbf Y}(c)
        }
        =
        1-\frac{0{,}38}{2}
        =
        0{,}81.
    \]

    Regresní přímka odstranila $81\,\%$ střední kvadratické chyby nejlepšího konstantního modelu.}
\end{frame}

\begin{frame}{Příklad}{Souvislost s~korelací}

    \visible<2->{Pro zadané soubory $\mathbf{X},\mathbf{Y}$ platí
    \[
        \sigma_\mathbf{X}
        =
        \sqrt{\frac15
        \sum_{i=1}^5
        (x_i-\average{\mathbf{X}})^2}
        =
        \sqrt{\frac{10}{5}}
        =
        \sqrt{2}
    \]
    a~stejně $\sigma_\mathbf{Y}=\sqrt{\frac15\sum_{i=1}^5(y_i-\average{\mathbf{Y}})^2}=\sqrt{2}$.\\[1mm]}

    \visible<3->{Kovariance je $\Cov(\mathbf X,\mathbf Y)=\frac15\cdot 9=1{,}8$. Proto
    \[
        \rho_{\mathbf X\mathbf Y}=
        \frac{
            \Cov(\mathbf X,\mathbf Y)
        }{
            \sigma_\mathbf{X}\sigma_\mathbf{Y}
        }=
        \frac{1{,}8}{\sqrt2\sqrt2}
        =
        0{,}9.
    \]}

    \visible<4->{Pro jednoduchou lineární regresi tedy skutečně platí
    \[
        \boxed{
            R^2=\rho_{\mathbf X\mathbf Y}^2=0{,}9^2=0{,}81
        }.
    \]}
\end{frame}

% =================================================================
\section{Omezení modelu}
% =================================================================

\begin{titleframeenv}{Omezení modelu}
    \obrazek[0.4\textheight]{images/omezeni-nelinearni-vztah.tex}
\end{titleframeenv}

\subsection{Kdy dává lineární regrese smysl?}

\begin{frame}{Kdy dává lineární regrese smysl?}

    \visible<2->{Před použitím lineární regrese očekáváme zejména:}

    \visible<3->{\begin{itemize}
        \item vztah mezi $\mathbf{X}$ a~$\mathbf{Y}$ je přibližně lineární,
        \visible<4->{\item pozorování nejsou vzájemně závislá,}
        \visible<5->{\item rezidua nevykazují systematický vzor,}
        \visible<6->{\item velikost reziduí se s~rostoucími hodnotami v~$\mathbf{X}$ výrazně nemění.}
    \end{itemize}}

    \visible<7->{\begin{tinfoframe}
        Tyto podmínky zde chápeme intuitivně: bodový graf a~rezidua by neměly ukazovat zjevný problém.
    \end{tinfoframe}}
\end{frame}

\subsection{Omezení}

\begin{frame}{Omezení}{Nelineární vztah}

    \visible<2->{Lineární regrese zachycuje \markdarkred{pouze lineární vztahy}.

    \begin{columns}[T]
        \begin{column}{0.45\textwidth}
            \begin{itemize}
                \item Mezi veličinami může být přesná závislost.
                \visible<3->{\item Přesto může být korelace nulová.}
                \visible<4->{\item Přímka pak vztah nevystihne.}
            \end{itemize}
        \end{column}

        \begin{column}{0.50\textwidth}\visible<5->{
    \obrazek[0.4\textheight]{images/omezeni-nelinearni-vztah.tex}
        }\end{column}
    \end{columns}}

    \visible<6->{\begin{talertframe}
        \begin{center}
            Nulová korelace neznamená, že mezi daty není žádný vztah.
        \end{center}
    \end{talertframe}}
\end{frame}


\begin{frame}{Omezení}{odlehlé hodnoty a~extrapolace}
    \begin{columns}[T]
        \begin{column}{0.47\textwidth}

            \visible<2->{\centering
            \textbf{Odlehlá hodnota}


            \obrazek[0.4\textheight]{images/omezeni-odlehle-hodnoty-a-extrapolace-1.tex}}
        \end{column}

        \begin{column}{0.47\textwidth}\visible<2->{

            \visible<3->{\centering
            \textbf{Extrapolace}


            \obrazek[0.4\textheight]{images/omezeni-odlehle-hodnoty-a-extrapolace-2.tex}}
        }\end{column}
    \end{columns}

    \visible<4->{\begin{talertframe}
        Jeden vzdálený bod může přímku výrazně změnit. Model také nemusí fungovat mimo rozsah naměřených dat.
    \end{talertframe}}
\end{frame}


\begin{frame}{Omezení}{Regrese není kauzalita}

    \visible<2->{Silný regresní vztah neznamená, že změna $\mathbf{X}$ způsobuje změnu $\mathbf{Y}$.}

    \visible<3->{\begin{talertbox}{Pozor}
        Lineární regrese popisuje vztah v~datech, ale sama nedokazuje příčinnou souvislost.
    \end{talertbox}}

    \visible<4->{\begin{itemize}
        \item Může existovat třetí proměnná, která ovlivňuje obě sledované veličiny.
        \visible<5->{\item Vztah může vzniknout náhodně.}
        \visible<6->{\item Závěr vždy závisí na kontextu a~smyslu dat.}
    \end{itemize}}

    \visible<7->{\begin{tinfoframe}
        Regrese odpovídá na otázku "jak spolu hodnoty souvisí v~datech", nikoli automaticky na~otázku "co je příčinou čeho".
    \end{tinfoframe}}
\end{frame}

% =================================================================
\section{Shrnutí}
% =================================================================

\begin{frame}{Shrnutí}

    \visible<2->{\begin{itemize}
        \item \textbf{Lineární regrese} hledá přímku
        \[
            f(x)=ax+b,
        \]
        která co nejlépe odhaduje hodnoty souboru $\mathbf Y$
        z~hodnot souboru $\mathbf X$.

        \visible<3->{\item Koeficienty určujeme \textbf{metodou nejmenších čtverců},
        tedy minimalizujeme
        \[
            \MSE_{\mathbf X,\mathbf Y}(f)
            =
            \frac1n\sum_{i=1}^n
            (y_i-f(x_i))^2.
        \]}

        \visible<4->{\item Nejlepší regresní přímka prochází bodem $(\average{\mathbf{X}},\average{\mathbf{Y}})$ a~její koeficienty jsou
        \[
            a=
            \frac{\Cov(\mathbf X,\mathbf Y)}
                 {\Var(\mathbf X)}
            =
            \rho_{\mathbf{X}\mathbf{Y}}\frac{\sigma_\mathbf{Y}}{\sigma_\mathbf{X}},
            \qquad
            b=\average{\mathbf{Y}}-a\average{\mathbf{X}}.
        \]}
    \end{itemize}}

\end{frame}

\begin{frame}{Shrnutí}
    \footnotesize

    \visible<2->{\begin{itemize}
        \item \textbf{Koeficient determinace} porovnává regresní přímku $f$ s~nejlepším konstantním modelem $c(x)=\average{\mathbf{Y}}$:
        \[
            R^2
            =
            1-
            \frac{
                \MSE_{\mathbf X,\mathbf Y}(f)
            }{
                \MSE_{\mathbf X,\mathbf Y}
                (x\mapsto\average{\mathbf{Y}})
            }.
        \]
    \end{itemize}}

    \visible<3->{\begin{timportantframe}
        Pro jednoduchou lineární regresi platí
        \[
            \boxed{R^2=\rho_{\mathbf{X}\mathbf{Y}}^2}.
        \]
        Čím je $R^2$ blíže jedné, tím větší část chyby konstantního modelu odstraníme použitím regresní přímky.
    \end{timportantframe}}

    \visible<4->{\begin{talertframe}
        Lineární regrese zachycuje pouze lineární vztah, je citlivá na~odlehlé hodnoty a~sama o~sobě nedokazuje příčinnou souvislost.
    \end{talertframe}}
    
\end{frame}

\zdrojeframe{russell-norvig-aima, james-isl-python}

\end{document}